\section{Sur la non-localité}
\subsection{Théorème EPR}
\subsubsection{Expérience de pensée}
Considérons deux particules de spin 1/2 constituant un système global dans un l’état
singulet de spin (de type produit tensoriel avec les deux états par ordre de particule)
\[
\ket{\Psi_s}=\frac{1}{\sqrt{2}}(\ket{+}\ket{-}-\ket{-}\ket{+})
\]
dit état intriqué. $\ket{+}(\ket{-})$ Représente un état propre de valeur propre
$+\hbar/2(-\hbar/2)$ de
l’observable projection du spin le long d’un axe (Oz). Ces deux particules sont émises dans
des directions opposées puis, une distance finie [non nulle] entre elles ayant été atteinte, l’on
mesure leur spin respectif grâce à des appareillages de type Stern \& Gerlach mesurant les
projections du spin selon un axe faisant un angle a avec (Oz) pour l’une et un angle b pour
l’autre. De par l’intrication, et négligeant l’effet d’évolution temporelle par simplicité (sans
impact sur la discussion), la mesure par exemple de la valeur $+\hbar/2$ le long de (Oz) pour
l’une des deux particules [projection de l’état singulet] donne immédiatement avec certitude
l’information à l’expérimentateur que l’état de l’autre particule – se trouvant pourtant à une
distance non nulle de lui – est l’état propre $\ket{-}$ de valeur propre $-\hbar/2$ (résultat assuré à
100\% d’une possible mesure). Il existe de même des corrélations entre les probabilités de
mesure des projections de spin pour les deux particules dans des directions caractérisées par
les angles a et b respectivement. Ces conclusions, souvent présentées comme un paradoxe
(EPR), violent, de manière plus précise, la loi relativiste (classique) d’une transmission
d’information ne pouvant être instantanée, bien qu’ici aucun système physique ne se soit
propagé à une vitesse (infinie) supérieure à celle de la lumière dans le vide.
\subsubsection{Non-séparabilité et causalité}
La MQ semble donc ne pas être « locale ». Même la « localité » non-relativiste [l’influence
physique décroît avec la distance ; notion purement spatiale elle] n’est pas respectée dans ce
cas de figure. Cette non-localité quantique provient d’une « non-séparabilité » : des
systèmes quantiques distincts ayant interagi peuvent ne plus posséder certaines propriétés
physiques propres, individuelles et faire alors partie d’un système quantique global (avec de
nouvelles propriétés) les réunissant. À noter que la non-séparabilité quantique est différente
de la non-séparation géographique.

Il existe aussi le principe de « causalité » (phénomène, appelé cause, précédant l’effet) en
sciences physiques, qui peut être interrogé en MQ du point de vue du non-déterminisme,
comme du point de vue temporel puisque certaines expériences paraissent montrer qu’un
3$^{ème}$ expérimentateur mesurant une polarisation de photon (dans un état superposé) ne peut
distinguer entre, le scénario où le photon a d’abord été manipulé par un 1$^{er}$ avant d’être
envoyé chez un 2$^{ème}$, et, le scénario inverse.
\subsubsection{Le théorème}
La conclusion qui précède peut être formulée de manière plus générale selon (Laloë) le
Théorème d’Einstein, Podolsky et Rosen (EPR) qui s’énonce comme suit, et définit le
contexte du « réalisme local ».
Si toutes les prédictions de mécanique quantique sont correctes (même pour des systèmes
constitués de plusieurs particules éloignées) et si la réalité physique peut être décrite dans
un cadre local (ou séparable), alors la mécanique quantique est nécessairement
incomplète : il existe dans la Nature des « éléments de réalité » qui sont laissés de côté par
cette théorie.

Nous synthétisons ce théorème dans le tableau ci-dessous, qui sera également utile pour la
suite.

2 interprétations de l'expérience EPR :

\begin{tabular}{rrc}
1)Localité \& éléments de réalité $\Rightarrow$ & $-2\leqslant<C>\leqslant+2$ (Bell)\\
2) Non-localité de MQ complète $\Rightarrow$&$-2\sqrt{2}\leqslant<C>\leqslant+2\sqrt{2}$\\
\end{tabular}
\subsection{Inégalités de Bell}
Appelons A(a,$\lambda$) et B(b,$\lambda$) les résultats [valeurs ± 1 en unités
 de $\hbar$/2] de la mesure des
projections respectives de spin des deux particules de l’expérience EPR (c.f. V.1.a). Ces
résultats dépendent bien des angles a,b de projection ainsi que possiblement d’éléments de
réalité notés $\lambda$ (pouvant dénoter un ensemble de paramètres) qui
 implémentent et doivent
expliquer les corrélations entre les mesures associées aux deux particules à distance finie
l’une de l’autre (c.f. V.1.c). Les éléments de réalité peuvent être liés
 à des variables cachées
(c.f. III.3). Considérons maintenant la combinaison (Eberhard) de résultats suivante,
\[
C+AB-AB'+A'B+A'B'=A(B-B')+A'(B+B')
\]

\subsection{Expériences d’A.Aspect}
\subsection{Ordre implicite de D.Bohm}
\[\]
\[\]
